Put \(A=2/p_c\) and \(L(R)=\log(2+R)\). The Kakeya bound of \ref{ms:cor-9-1}, with power \(2A\), gives the bound \(C L(1/\rho)^{2A}\) in \ref{ms:nikodym-distal} for the supremum \(D_\rho\) over neighborhoods of the ray segments at distances between \(1/2\) and \(1\) from \(x\).

Take any unit segment through \(x\). Its points within \(2\delta\) of \(x\), with their \(\delta\)-neighborhood, lie in \(B(x,3\delta)\). On each of the two rays the rest lies in the union of the intervals of distances \([s/2,s]\) with \(s=2^{-j}\), \(j\ge0\), \(s\ge\delta\). For such \(s\), a \(\delta\)-neighborhood of the ray interval has volume \(\lesssim_n s\delta^{n-1}\), as \(\delta\le s\). The full tube volume is comparable to \(\delta^{n-1}\); thus the ball and ray intervals contribute at most a constant times \(\delta\) and \(s\), respectively, times their individual maximal averages. The ball average has \(L^n\) operator norm one by Young's inequality. The maximal ray average at scale \(s\) is conjugate by dilation to \(D_{\delta/s}\), so its norm is bounded by \(C L(s/\delta)^{2A}\). Consequently
\[
 \|N_\delta\|_{L^n\to L^n}
 \lesssim \delta+\sum_{\substack{s=2^{-j},\ j\ge0\\ s\ge\delta}}sL(s/\delta)^{2A}
 \lesssim L(1/\delta)^{2A}.
\]
